Across the 216 model--asset pairs, ACI with $\gamma$ selected by first-half validation --- the grid value minimising the absolute gap between held-out online coverage and the 0.01 nominal rate on the first half of the calibration sample --- attains a Green-zone share of 90.3\% at a mean violation rate of 0.006, against 94.9\% at 0.010 for the 250-day rolling estimator (target $0.01$); the two are within 4.6 percentage points on Green-zone coverage, the rolling estimator marginally ahead. The separation is operational, and it runs against ACI on both axes it is usually credited with: ACI over-covers ($\hat\pi = 0.006$) and pays for it in capital, posting a mean $|\VaR|$ of 0.0750 against 0.0443 for the rolling estimator (a factor of 1.7), while also producing the rougher corrected path (day-over-day volatility 0.00992 versus 0.00579, a factor of 1.7). No step size in the grid overturns the capital ordering: the lightest-capital setting still posts a mean $|\VaR|$ of 0.0646 against 0.0443 for the rolling estimator (Table~\ref{tab:aci_gamma_sensitivity}). On this evidence the rolling estimator dominates ACI at the Basel/capital operating point --- comparable zone coverage at materially lower capital and a smoother reported path --- and remains the headline procedure, backed by its transparent finite-window coverage bound (Remark~\ref{rem:rolling}). ACI nonetheless remains attractive when the loss function penalises zone placement over capital and a fully online, distribution-free update is required.
